\documentclass[11pt]{article}
\usepackage[margin=1in]{geometry}
\usepackage{amsmath,amssymb,amsthm}
\usepackage[colorlinks=true,urlcolor=blue,linkcolor=blue,citecolor=blue]{hyperref}
\newtheorem{theorem}{Theorem}
\newtheorem{lemma}{Lemma}
\newcommand{\Med}{\operatorname{Med}}
\newcommand{\Prob}{\mathbb P}
\newcommand{\E}{\mathbb E}
\title{The McKay Type I median: strict monotonicity\\
and correction of a conjectured lower bound}
\author{Principia Math in Progress}
\date{12 September 2026}
\begin{document}
\maketitle

\paragraph{Task 1: Resolve MathDB 345411.}
Gaunt--Merkle's Conjecture 3.6 \cite{GM} concerns the McKay Type I law
with \(m>-\tfrac12\), \(c>1\), and
\(\phi=bc/(c^2-1)>0\) fixed. We prove the intended strict monotonicity
in \(c\) and the sharp gamma-median envelopes for every such parameter.
The full printed conjecture is nevertheless false: its exponential
lower bound is too large. An exact counterexample and a corrected
bound appear below. The printed map with variable \(\phi\), while
holding \(\phi\) fixed and varying \(c\), is treated as a notational
slip; it is not the basis of the refutation.

Write \(\Gamma(r,1)\) for the gamma law with shape \(r\) and unit rate,
and let \(q(r)=\Med(\Gamma(r,1))\).

\begin{theorem}\label{thm:main}
For fixed \(m>-\tfrac12\) and \(\phi>0\), the function
\(c\mapsto\Med(Z_{m,c,\phi})\) is strictly increasing on \((1,\infty)\).
Its endpoint limits give the sharp bounds
\begin{equation}\label{eq:envelopes}
 2\phi q(m+\tfrac12)<\Med(Z_{m,c,\phi})<\phi q(2m+1).
\end{equation}
For every \(m>-\tfrac12\) one consequently has
\begin{align}
 (2m+1)\phi\exp\!\left(-\frac{2\log2}{2m+1}\right)
 &<\Med(Z_{m,c,\phi}),                                      \label{eq:correct}\\
 \Med(Z_{m,c,\phi})
 &<(2m+1)\phi\exp\!\left(-\frac{1}{3(2m+1)}\right) \notag\\
 &<\left(2m+\frac23+\frac{1}{18(2m+1)}\right)\phi,             \label{eq:upper}\\
 (2m+\tfrac13)\phi&<\Med(Z_{m,c,\phi}).                       \label{eq:linear}
\end{align}
For \(m\geq0\) there is also the strict upper bound
\begin{equation}\label{eq:linearupper}
 \Med(Z_{m,c,\phi})<(2m+\log2)\phi.
\end{equation}
In contrast, at
\[
 m=0,\qquad \phi=1,\qquad c=\frac{1001}{1000},
 \qquad b=\frac{2001}{1001000},
\]
the median is less than \(12106/25025<1/2\), whereas the conjecture's
printed exponential lower bound equals \(1/2\).
\end{theorem}

\paragraph{Reduction to two gamma variables.}
Set \(a=m+\tfrac12>0\) and \(t=1/c\in(0,1)\). If \(X,Y\) are independent
\(\Gamma(a,1)\) variables, then
\begin{equation}\label{eq:sum}
 Z_{m,c,\phi}/\phi\ \stackrel{d}{=}\ 
 S_{a,t}:=(1+t)X+(1-t)Y .
\end{equation}
This is the representation underlying \cite[eq.\ (3.44)]{GM}.
Although its surrounding text states a narrower range, the gamma
convolution, or its Laplace transform below and the Bessel density,
establishes \eqref{eq:sum} for all \(a>0\).
Let \(F_{a,t},f_{a,t}\), and \(M_a(t)\) denote its distribution function,
density, and unique median. Its density is positive for \(x>0\).

\paragraph{A derivative identity.}
For \(s>0\),
\[
 L_{a,t}(s):=\E e^{-sS_{a,t}}
       =[1+2s+(1-t^2)s^2]^{-a}.
\]
The Laplace transform in \(x\) of \(F_{a,t}(x)\) is \(L_{a,t}(s)/s\).
Differentiation gives
\[
 \partial_t\{L_{a,t}(s)/s\}=2at\,sL_{a+1,t}(s).
\]
Since \(f_{a+1,t}(x)=O(x^{2a+1})\) at zero, its value there is zero.
Laplace uniqueness therefore yields
\begin{equation}\label{eq:identity}
 \partial_t F_{a,t}(x)=2at\,f'_{a+1,t}(x),\qquad x>0,\quad 0<t<1.
\end{equation}
These operations can be justified directly with the gamma convolution:
on a compact subinterval of \((-1,1)\), its parameter derivatives
are dominated by integrable functions with exponential tails.
The distribution functions and the displayed derivatives are
continuous in \(x>0\), so Laplace uniqueness gives pointwise equality.

\begin{lemma}\label{lem:increasing}
For \(a>0\), \(0<t<1\), and \(0<x\leq2a\), one has
\(f'_{a+1,t}(x)>0\).
\end{lemma}
\begin{proof}
Put \(c=1/t\), \(b=(1-t^2)/t\), and \(\nu=a+\tfrac12\).
The gamma convolution gives, with a positive constant \(C\),
\[
 f_{a+1,t}(x)=C x^\nu e^{-cx/b}I_\nu(x/b).
\]
We recall, and briefly prove, the Bessel-ratio bound
\begin{equation}\label{eq:ratio}
 \frac{I_\nu(u)}{I_{\nu-1}(u)}
 < R(u):=\frac{u}{a+\sqrt{a^2+u^2}},\qquad u>0 .
\end{equation}
The ratio \(Q=I_\nu/I_{\nu-1}\) satisfies
\[
 Q'=1-\frac{2a}{u}Q-Q^2.
\]
Here \(Q(u)\sim u/(2a+1)\), \(R(u)\sim u/(2a)\) at zero,
\(R'>0\), and \(1-2aR/u-R^2=0\).
At a first contact from below, \(Q'=0<R'\), which is impossible.
This proves \eqref{eq:ratio}, also recorded in \cite[eq.\ (2.10)]{GM}.

For \(u=x/b\leq2ac/(c^2-1)\), algebra gives
\(R(u)\leq1/c\). Indeed \(R\) is increasing and
\(R(2ac/(c^2-1))=1/c\).
Using \((x^\nu I_\nu(x/b))'=b^{-1}x^\nu I_{\nu-1}(x/b)\),
we obtain
\[
 \frac{f'_{a+1,t}(x)}{f_{a+1,t}(x)}
 =\frac1b\left(\frac{I_{\nu-1}(x/b)}{I_\nu(x/b)}-c\right)>0.
\]
\end{proof}

\paragraph{The median lies below the mean.}
First, \(q(r)<r\) for every \(r>0\). Here is a short direct proof.
For \(U\sim\Gamma(r,1)\), the density of \(V=\log(U/r)\) is
\[
 g(v)=\frac{r^r}{\Gamma(r)}\exp\{r(v-e^v)\}.
\]
For \(v>0\),
\[
 \frac{g(-v)}{g(v)}
 =\exp\{r(e^v-e^{-v}-2v)\}>1.
\]
Integration gives \(\Prob(U<r)>1/2\).
At \(t=0\), \(S_{a,0}\sim\Gamma(2a,1)\). Lemma
\ref{lem:increasing} and \eqref{eq:identity}, evaluated at \(x=2a\),
now imply
\[
 F_{a,t}(2a)>F_{a,0}(2a)>\frac12,\qquad 0<t<1.
\]
Continuity at \(t=0\) follows either from the convolution or from the
coupling \eqref{eq:sum}. Thus \(0<M_a(t)<2a\).

\paragraph{Strict monotonicity and the endpoint envelopes.}
Differentiate \(F_{a,t}(M_a(t))=1/2\), using the positive density
and the implicit function theorem:
\begin{equation}\label{eq:median}
 M_a'(t)
 =-\frac{2at\,f'_{a+1,t}(M_a(t))}{f_{a,t}(M_a(t))}<0.
\end{equation}
Since \(t=1/c\), this proves the asserted strict increase in \(c\).
The same coupling converges almost surely to \(X+Y\) as \(t\to0\)
and to \(2X\) as \(t\to1\). Both limiting laws have continuous,
strictly increasing distribution functions on their support, so their
medians converge:
\[
 M_a(0)=q(2a),\qquad M_a(1)=2q(a).
\]
Strict monotonicity in the interior gives
\(2q(a)<M_a(t)<q(2a)\), proving \eqref{eq:envelopes}.
The argument also resolves the equivalent two-weight formulation
in \cite[Conjecture 3.2]{GM}; no separate problem credit is taken.

\paragraph{The elementary bounds and the incorrect exponent.}
The established gamma-median estimates \cite{BP,GM} give
\[
 r e^{-(\log2)/r}<q(r)<r e^{-1/(3r)},\qquad
 q(r)>r-\tfrac13 \quad(r>0),
\]
and \(q(r)\leq r-1+\log2\) for \(r\geq1\), with equality only at \(r=1\).
Apply the lower estimates to \(2q(a)\), and the upper estimates to
\(q(2a)\). This proves \eqref{eq:correct}--\eqref{eq:linearupper};
the last step of \eqref{eq:upper} uses \(e^{-u}<1-u+u^2/2\), \(u>0\).
Also \(\log2>2/3\), so \eqref{eq:linear} implies the conjecture's
weaker linear lower bound \((2m+1-\log2)\phi\).
The lower endpoint has shape \(a=(2m+1)/2\), which explains the
factor \(2\) required in the exponent of \eqref{eq:correct}.

\paragraph{An exact refutation of the printed exponential bound.}
Use the parameters in Theorem \ref{thm:main}. Then
\[
 Z=\frac{2001}{1001}X+\frac1{1001}Y,\qquad
 X,Y\ \hbox{independent }\Gamma(\tfrac12,1).
\]
On the event \(X\leq6/25,\ Y\leq4\), one has
\[
 Z\leq \frac{2001}{1001}\frac6{25}+\frac4{1001}
      =\frac{12106}{25025}<\frac12 .
\]
This event has probability
\(\operatorname{erf}(\sqrt6/5)\operatorname{erf}(2)\).
Integration of \(e^{-u^2}>1-u^2\) gives
\[
 \operatorname{erf}(\sqrt6/5)>
 \frac{46\sqrt6}{125\sqrt\pi}.
\]
Moreover,
\[
 \operatorname{erfc}(2)
 <\frac{e^{-4}}{2\sqrt\pi}<\frac1{100},
 \qquad
 e^4>\sum_{j=0}^7\frac{4^j}{j!}=\frac{16319}{315}>50 .
\]
For the first inequality, bound
\(\int_2^\infty e^{-u^2}\,du\) by
\(\int_2^\infty(u/2)e^{-u^2}\,du\).
Using \(\pi<22/7\), the square of the resulting probability lower
bound satisfies
\[
 \left(\frac{99}{100}\frac{46\sqrt6}{125\sqrt\pi}\right)^2
 >\frac{9898119}{39062500}
 =\frac14+\frac{66247}{19531250}>\frac14.
\]
Thus \(\Prob(Z\leq12106/25025)>1/2\), proving the strict median bound.
At \(m=0,\phi=1\), the conjectured lower expression
\((2m+1)\phi e^{-(\log2)/(2m+1)}\) is \(1/2\), a contradiction.
The source's own Table 2 already reports a numerical value below
\(1/2\) at \(m=0,c=1.01\); that numerical inconsistency is not new.
The argument here supplies an exact certificate and a proof of the
central monotonicity assertion.

\paragraph{Verification and status.}
The companion verifier checks the transform identity, the rational
counterexample arithmetic, the Bessel comparison algebra, and
independent numerical distribution calculations. Numerical checks
are supplementary to the proof. The dated literature audit records
the current open MathDB entry, the retained conjecture in the current
author manuscript and indexed published text, and searches for
subsequent public claims. No matching full claim was located.

\paragraph{Task 2: Determine the quantile range of monotonicity.}
For each \(a>0\), determine exactly which probabilities \(p\) make the
\(p\)-quantile of \(S_{a,t}\) decrease throughout \(0<t<1\), and locate
any changes of direction for the remaining quantiles.

\paragraph{Task 3: Compute the endpoint expansions of the median.}
Derive the first nonzero correction terms to \(M_a(0)=q(2a)\) and
\(M_a(1)=2q(a)\), including uniform error bounds as \(a\) varies over
compact subsets of \((0,\infty)\).

\paragraph{Task 4: Extend the median comparison to more gamma summands.}
For independent identically distributed gamma variables, determine
whether the median of their weighted sum is Schur-concave in positive
weights with fixed total, and characterize strictness and boundary cases.

\paragraph{Task 5: Resolve the corresponding gamma-difference median problem.}
Develop an analogue of \eqref{eq:identity} for one negative weight and
use it to investigate the variance-gamma monotonicity and limiting
median in \cite[Conjectures 3.3 and 3.5]{GM}.

\begin{thebibliography}{9}
\bibitem{GM}
R.\ E.\ Gaunt and M.\ Merkle,
On bounds for the mode and median of the generalized hyperbolic and
related distributions,
\emph{J.\ Math.\ Anal.\ Appl.} \textbf{493} (2021), 124508.
\href{https://doi.org/10.1016/j.jmaa.2020.124508}{doi:10.1016/j.jmaa.2020.124508}.
Current author text:
\href{https://arxiv.org/abs/2002.01884}{arXiv:2002.01884v4}.
\bibitem{BP}
C.\ Berg and H.\ L.\ Pedersen,
The Chen--Rubin conjecture in a continuous setting,
\emph{Methods Appl.\ Anal.} \textbf{13} (2006), 63--88.
\end{thebibliography}
\end{document}
